\documentclass[../DoAn.tex]{subfiles}
\begin{document}

Chương 2 đã trình bày các kiến thức nền tảng. Chương này mô tả quá trình \textit{phân tích yêu cầu} và \textit{thiết kế hệ thống}: xác định các chức năng cần có, vẽ kiến trúc tổng thể, thiết kế cơ sở dữ liệu, các API, và quy trình phân tích CV bằng phương pháp lai.

\section{Phân tích yêu cầu}

\subsection{Yêu cầu chức năng}

Dựa trên bài toán đã định nghĩa ở Chương 1, hệ thống cần đáp ứng các yêu cầu chức năng liệt kê trong Bảng \ref{tab:func_req}.

\begin{table}[H]
\centering
\caption{Danh sách yêu cầu chức năng}
\label{tab:func_req}
\resizebox{\textwidth}{!}{%
\begin{tabular}{|c|l|p{8cm}|}
\hline
\textbf{Mã} & \textbf{Tên yêu cầu} & \textbf{Mô tả} \\ \hline
FR01 & Đăng ký tài khoản & Người dùng đăng ký bằng email và mật khẩu \\ \hline
FR02 & Đăng nhập đa phương thức & Đăng nhập bằng email/mật khẩu, Google OAuth, hoặc LinkedIn OAuth \\ \hline
FR03 & Quản lý phiên đăng nhập & Duy trì phiên đăng nhập bằng JWT, hỗ trợ đăng xuất \\ \hline
FR04 & Hỗ trợ phiên khách & Cho phép sử dụng trình soạn thảo CV mà không cần đăng nhập \\ \hline
FR05 & Tải lên CV & Upload file PDF hoặc DOCX (tối đa 10 MB) \\ \hline
FR06 & Trích xuất văn bản & Tự động trích xuất văn bản từ tệp PDF và DOCX \\ \hline
FR07 & Phân tích CV bằng AI & Phân tích nội dung CV, trích xuất thông tin có cấu trúc \\ \hline
FR08 & Chấm điểm tin cậy & Tính điểm tin cậy cho kết quả phân tích \\ \hline
FR09 & Nhận diện ngôn ngữ & Tự động phát hiện CV viết bằng tiếng Việt hay tiếng Anh \\ \hline
FR10 & Chỉnh sửa theo phần & Editor cho phép sửa từng section: liên hệ, tóm tắt, kinh nghiệm, học vấn, kỹ năng \\ \hline
FR11 & Kéo thả sắp xếp & Cho phép sắp xếp thứ tự các mục bằng thao tác kéo thả \\ \hline
FR12 & Tự động lưu & Auto-save dữ liệu CV định kỳ \\ \hline
FR13 & Xem trước real-time & Preview CV định dạng A4 cập nhật theo thời gian thực \\ \hline
FR14 & Xuất PDF & Xuất CV ra định dạng PDF \\ \hline
FR15 & Xuất DOCX & Xuất CV ra định dạng Microsoft Word \\ \hline
FR16 & Quản lý workspace & Dashboard hiển thị, tạo mới, sửa, xoá CV \\ \hline
FR17 & Hỗ trợ song ngữ & Giao diện hỗ trợ cả tiếng Việt và tiếng Anh \\ \hline
FR18 & So khớp CV-JD & So sánh CV với mô tả công việc bằng ba phương pháp matching \\ \hline
\end{tabular}%
}
\end{table}

\subsection{Yêu cầu phi chức năng}

Bảng \ref{tab:nonfunc_req} liệt kê các yêu cầu phi chức năng quan trọng.

\begin{table}[H]
\centering
\caption{Danh sách yêu cầu phi chức năng}
\label{tab:nonfunc_req}
\begin{tabular}{|c|l|p{8cm}|}
\hline
\textbf{Mã} & \textbf{Tên yêu cầu} & \textbf{Mô tả} \\ \hline
NFR01 & Hiệu năng & Thời gian phân tích CV dưới 15 giây \\ \hline
NFR02 & Khả năng mở rộng & Kiến trúc module hỗ trợ thêm tính năng mới \\ \hline
NFR03 & Tương thích trình duyệt & Hoạt động trên Chrome, Firefox, Safari, Edge \\ \hline
\end{tabular}
\end{table}

\subsection{Biểu đồ Use Case}

Hệ thống có ba tác nhân chính: Khách (chưa đăng nhập), Người dùng (đã đăng nhập), và Quản trị viên. Hình \ref{fig:usecase} mô tả tổng thể các trường hợp sử dụng.

\begin{figure}[H]
\centering
\begin{tikzpicture}[
  every node/.style={font=\small},
  uc/.style={ellipse, draw, thick, minimum width=3.8cm, minimum height=1.0cm, align=center, font=\footnotesize, fill=blue!8},
  actor/.style={font=\footnotesize\bfseries, align=center},
  sysbox/.style={rectangle, draw, very thick, dashed, inner sep=0.6cm},
  conn/.style={thick, gray!70}
]
\foreach \name/\y/\label in {guest/4.0/Khách, user/-1.1/Người dùng, admin/-5.9/Quản trị viên} {
  \node[circle, draw, minimum size=0.45cm, inner sep=0pt] (\name h) at (-5, \y+0.4) {};
  \draw[thick] (-5, \y+0.18) -- (-5, \y-0.6);
  \draw[thick] (-5.4, \y-0.2) -- (-4.6, \y-0.2);
  \draw[thick] (-5, \y-0.6) -- (-5.35, \y-1.05);
  \draw[thick] (-5, \y-0.6) -- (-4.65, \y-1.05);
  \node[actor, anchor=north] at (-5, \y-1.1) {\label};
  \coordinate (\name) at (-4.5, \y-0.2);
}

\node[uc] (uc1) at (2, 5.4)  {Xem trang chủ};
\node[uc] (uc2) at (2, 4.0)  {Soạn thảo CV};
\node[uc] (uc3) at (2, 2.6)  {Xuất CV\\(PDF/DOCX)};
\node[uc] (uc4) at (2, 1.0)  {Đăng nhập / Đăng ký};
\node[uc] (uc5) at (2, -0.4) {Tải lên CV};
\node[uc] (uc6) at (2, -1.8) {Phân tích CV\\bằng LLM};
\node[uc] (uc7) at (2, -3.2) {Quản lý CV\\(workspace)};
\node[uc] (uc8) at (2, -5.2) {Xem thống kê};
\node[uc] (uc9) at (2, -6.6) {Quản lý người dùng};

\begin{scope}[on background layer]
  \node[sysbox, fit=(uc1)(uc2)(uc3)(uc4)(uc5)(uc6)(uc7)(uc8)(uc9), label={[font=\bfseries\footnotesize]above:Hệ thống}] {};
\end{scope}

\draw[conn] (guest) -- (uc1.west);
\draw[conn] (guest) -- (uc2.west);
\draw[conn] (guest) -- (uc3.west);
\draw[conn] (user) -- (uc4.west);
\draw[conn] (user) -- (uc5.west);
\draw[conn] (user) -- (uc6.west);
\draw[conn] (user) -- (uc7.west);
\draw[conn] (admin) -- (uc8.west);
\draw[conn] (admin) -- (uc9.west);
\end{tikzpicture}
\caption{Biểu đồ Use Case tổng thể của hệ thống}
\label{fig:usecase}
\end{figure}

Khách có thể truy cập trang chủ, soạn thảo CV từ template, và xuất CV mà không cần đăng nhập. Người dùng đã đăng nhập có thêm các chức năng quản lý workspace, lưu CV vào cơ sở dữ liệu, tải lên CV để phân tích. Quản trị viên có quyền truy cập dashboard quản trị và xem thống kê hệ thống.

\section{Kiến trúc tổng thể hệ thống}

Hệ thống được thiết kế theo kiến trúc \textit{ba lớp} (three-tier) gồm: lớp trình bày, lớp logic nghiệp vụ, và lớp dữ liệu. Toàn bộ hệ thống chạy trên nền tảng Next.js cho phép tích hợp cả frontend và backend API trong cùng một dự án. Hình \ref{fig:architecture} mô tả tổ chức ba lớp.

\begin{figure}[H]
\centering
\begin{tikzpicture}[
  every node/.style={font=\small},
  layer/.style={rectangle, rounded corners, draw, very thick, minimum width=13cm, minimum height=2.4cm, align=center},
  comp/.style={rectangle, rounded corners=2pt, draw, minimum width=2.5cm, minimum height=1.0cm, align=center, font=\footnotesize, fill=white},
  arr/.style={{Stealth[length=2.5mm]}-{Stealth[length=2.5mm]}, thick}
]
\node[layer, fill=blue!10] (pres) {};
\node[anchor=north, font=\bfseries, yshift=-0.15cm] at (pres.north) {Lớp trình bày};
\node[comp] at ([xshift=-4.5cm, yshift=-0.45cm]pres.center) {Landing\\Page};
\node[comp] at ([xshift=-1.5cm, yshift=-0.45cm]pres.center) {CV Upload\\\& Editor};
\node[comp] at ([xshift=1.5cm, yshift=-0.45cm]pres.center) {Workspace};
\node[comp] at ([xshift=4.5cm, yshift=-0.45cm]pres.center) {Admin\\Dashboard};

\node[layer, fill=orange!12, below=0.5cm of pres] (logic) {};
\node[anchor=north, font=\bfseries, yshift=-0.15cm] at (logic.north) {Lớp logic nghiệp vụ};
\node[comp] at ([xshift=-4.5cm, yshift=-0.45cm]logic.center) {API\\Routes};
\node[comp] at ([xshift=-1.5cm, yshift=-0.45cm]logic.center) {Xác thực\\(JWT, OAuth)};
\node[comp] at ([xshift=1.5cm, yshift=-0.45cm]logic.center) {Hybrid CV\\Parser};
\node[comp] at ([xshift=4.5cm, yshift=-0.45cm]logic.center) {Matching\\Engine};

\node[layer, fill=green!12, below=0.5cm of logic] (data) {};
\node[anchor=north, font=\bfseries, yshift=-0.15cm] at (data.north) {Lớp dữ liệu};
\node[comp] at ([xshift=-4.5cm, yshift=-0.45cm]data.center) {Supabase\\PostgreSQL};
\node[comp] at ([xshift=-1.5cm, yshift=-0.45cm]data.center) {pgvector\\(1536-dim)};
\node[comp] at ([xshift=1.5cm, yshift=-0.45cm]data.center) {Vercel Blob\\(File storage)};
\node[comp] at ([xshift=4.5cm, yshift=-0.45cm]data.center) {OpenAI API};

\draw[arr] (pres) -- (logic);
\draw[arr] (logic) -- (data);
\end{tikzpicture}
\caption{Kiến trúc tổng thể hệ thống ba lớp}
\label{fig:architecture}
\end{figure}

\textbf{Lớp trình bày} là các trang giao diện viết bằng React/Next.js: trang chủ, trang đăng nhập, trang tải lên CV, trình soạn thảo, workspace quản lý CV, và admin dashboard.

\textbf{Lớp logic nghiệp vụ} gồm các API endpoint Next.js xử lý: xác thực (JWT, OAuth), hybrid CV parser (regex + LLM), matching engine (TF-IDF, Embedding, LLM).

\textbf{Lớp dữ liệu} dùng Supabase PostgreSQL kèm tiện ích mở rộng pgvector để lưu vector 1536 chiều. Tệp CV được lưu trên Vercel Blob. Các lời gọi tới OpenAI API (cho LLM và embedding) cũng được tính trong lớp này.

\section{Thiết kế Mô hình thực thể và Cơ sở dữ liệu}

\subsection{Mô hình thực thể (Domain Model)}

Ở mức ứng dụng, hệ thống mô hình hoá các thực thể nghiệp vụ dưới dạng các lớp có quan hệ rõ ràng. Hình \ref{fig:domain_model} trình bày biểu đồ lớp của các thực thể chính.

\begin{figure}[H]
\centering
\begin{tikzpicture}[
  cls/.style={rectangle, draw, very thick, fill=white, inner sep=0pt},
  rel/.style={thick, -{Stealth[length=2mm]}},
  rellabel/.style={font=\scriptsize, fill=white, inner sep=1pt}
]
\node[cls] (User) at (0, 0) {%
\begin{tabular}{|p{3.6cm}|}
\hline
\rowcolor{blue!15}\centering\bfseries User \tabularnewline\hline
\scriptsize\,-- id: UUID \tabularnewline
\scriptsize\,-- email: string \tabularnewline
\scriptsize\,-- passwordHash: string? \tabularnewline
\scriptsize\,-- oauthProvider: string? \tabularnewline
\scriptsize\,-- emailVerified: bool \tabularnewline
\scriptsize\,-- createdAt: datetime \tabularnewline\hline
\end{tabular}};

\node[cls] (CV) at (8, 0) {%
\begin{tabular}{|p{3.6cm}|}
\hline
\rowcolor{blue!15}\centering\bfseries CV \tabularnewline\hline
\scriptsize\,-- id: UUID \tabularnewline
\scriptsize\,-- userId: UUID? \tabularnewline
\scriptsize\,-- title: string \tabularnewline
\scriptsize\,-- cvData: JSONB \tabularnewline
\scriptsize\,-- jobDescription: string? \tabularnewline
\scriptsize\,-- possibilityScore: int \tabularnewline
\scriptsize\,-- createdAt: datetime \tabularnewline\hline
\end{tabular}};

\node[cls] (Exp) at (0, -5.3) {%
\begin{tabular}{|p{3.6cm}|}
\hline
\rowcolor{blue!15}\centering\bfseries Experience \tabularnewline\hline
\scriptsize\,-- title: string \tabularnewline
\scriptsize\,-- company: string \tabularnewline
\scriptsize\,-- startDate: string \tabularnewline
\scriptsize\,-- endDate: string? \tabularnewline
\scriptsize\,-- isCurrent: bool \tabularnewline
\scriptsize\,-- bullets: string[] \tabularnewline\hline
\end{tabular}};

\node[cls] (Section) at (8, -5.3) {%
\begin{tabular}{|p{3.6cm}|}
\hline
\rowcolor{blue!15}\centering\bfseries Section \tabularnewline\hline
\scriptsize\,-- id: string \tabularnewline
\scriptsize\,-- type: enum (contact, \tabularnewline
\scriptsize\quad summary, experience, \tabularnewline
\scriptsize\quad education, skills, custom) \tabularnewline
\scriptsize\,-- order: int \tabularnewline
\scriptsize\,-- content: JSONB \tabularnewline\hline
\end{tabular}};

\node[cls] (Guest) at (0, -10) {%
\begin{tabular}{|p{3.6cm}|}
\hline
\rowcolor{blue!15}\centering\bfseries GuestSession \tabularnewline\hline
\scriptsize\,-- sessionId: UUID \tabularnewline
\scriptsize\,-- cvDataDraft: JSONB \tabularnewline
\scriptsize\,-- expiresAt: datetime \tabularnewline\hline
\end{tabular}};

\draw[rel] (User.east) -- node[rellabel, above, pos=0.08]{1} node[rellabel, above, pos=0.92]{*} (CV.west);
\draw[rel] (CV.south) -- node[rellabel, right, pos=0.1]{1} node[rellabel, right, pos=0.9]{*} (Section.north);
\draw[rel] (Section.west) -- node[rellabel, above, pos=0.08]{1} node[rellabel, above, pos=0.92]{*} (Exp.east);
\draw[rel, dashed] (Guest.west) -- ++(-1.2,0) |- node[rellabel, pos=0.25, left]{nâng cấp} (User.west);
\end{tikzpicture}
\caption{Mô hình thực thể (Domain Model) các đối tượng nghiệp vụ chính}
\label{fig:domain_model}
\end{figure}

Mỗi \texttt{User} có thể sở hữu nhiều \texttt{CV}. Mỗi \texttt{CV} gồm nhiều \texttt{Section} có thứ tự (kiểu enum: liên hệ, tóm tắt, kinh nghiệm, học vấn, kỹ năng, tuỳ chỉnh); trong đó section kinh nghiệm chứa danh sách \texttt{Experience}. Thực thể \texttt{GuestSession} cho phép khách dùng trình soạn thảo mà không cần đăng nhập; dữ liệu nháp có thể được nâng cấp thành \texttt{CV} thuộc \texttt{User} mới khi khách đăng ký.

\subsection{Sơ đồ quan hệ thực thể (ERD)}

Ở mức cơ sở dữ liệu, hệ thống dùng sáu bảng chính. Hình \ref{fig:erd} mô tả quan hệ giữa chúng.

\begin{figure}[H]
\centering
\begin{tikzpicture}[
  every node/.style={font=\scriptsize},
  header/.style={rectangle, draw, very thick, fill=blue!15, minimum width=4cm, minimum height=0.5cm, font=\footnotesize\bfseries, align=center, anchor=north west},
  fields/.style={rectangle, draw, very thick, fill=white, minimum width=4cm, anchor=north west, align=left, text width=3.8cm, inner sep=3pt, font=\scriptsize},
  arr/.style={-{Stealth[length=2mm]}, thick, blue!70!black},
  one/.style={font=\scriptsize\bfseries, fill=white, inner sep=1pt},
  many/.style={font=\scriptsize\bfseries, fill=white, inner sep=1pt}
]
\node[header] (uh) at (0, 4) {users};
\node[fields, below=0pt of uh.south west] (uf) {
- id (PK) UUID\\
- email\\
- password\_hash\\
- oauth\_provider\\
- created\_at};

\node[header] (wh) at (6, 4) {cv\_workflow};
\node[fields, below=0pt of wh.south west] (wf) {
- id (PK) TEXT\\
- user\_id (FK $\to$ users)\\
- title\\
- cv\_data JSONB\\
- status\\
- cv\_embedding vector(1536)};

\node[header] (dh) at (12, 4) {cv\_drafts};
\node[fields, below=0pt of dh.south west] (df) {
- id (PK) UUID\\
- user\_id (FK $\to$ users)\\
- file\_id\\
- jd\_text\\
- analysis\_id};

\node[header] (ch) at (0, -2.0) {jd\_targets};
\node[fields, below=0pt of ch.south west] (cf) {
- id (PK) UUID\\
- cv\_id (FK $\to$ cv\_workflow)\\
- label\\
- jd\_text\\
- jd\_embedding vector(1536)};

\node[header] (sh) at (6, -2.0) {match\_runs};
\node[fields, below=0pt of sh.south west] (sf) {
- id (PK) UUID\\
- cv\_id\\
- jd\_target\_id\\
- method (tfidf/embed/llm)\\
- score, latency, cost};

\node[header] (ah) at (12, -2.0) {audit\_logs};
\node[fields, below=0pt of ah.south west] (af) {
- id (PK) UUID\\
- user\_id (FK $\to$ users)\\
- action\\
- timestamp};

\draw[arr] (uf.east |- uh.south) -- node[one, pos=0.1]{1} node[many, pos=0.9]{*} (wf.west |- wh.south);
\draw[arr] (wf.east |- wh.south) -- node[one, pos=0.1]{1} node[many, pos=0.9]{*} (df.west |- dh.south);
\draw[arr] ($(uf.south)+(0,0)$) -- ++(0,-0.5) -| node[many, pos=0.95]{*} (cf.north);
\draw[arr] (cf.east |- ch.south) -- node[one, pos=0.1]{1} node[many, pos=0.9]{*} (sf.west |- sh.south);
\draw[arr] ($(uf.south)+(1.0,0)$) -- ++(0,-0.7) -| node[many, pos=0.95]{*} (af.north);
\end{tikzpicture}
\caption{Sơ đồ quan hệ thực thể (ERD) của hệ thống}
\label{fig:erd}
\end{figure}

\subsection{Mô tả các bảng dữ liệu}

\textbf{Bảng \texttt{users}} lưu thông tin tài khoản: họ tên, email, mật khẩu đã băm (bcrypt), trạng thái xác minh, và thông tin OAuth. Ràng buộc: nếu đăng nhập bằng OAuth thì không cần mật khẩu, ngược lại phải có mật khẩu đã băm.

\textbf{Bảng \texttt{cv\_workflow}} là bảng trung tâm lưu trữ CV đầy đủ. Trường \texttt{cv\_data} kiểu JSONB lưu toàn bộ cấu trúc CV (các section và nội dung). Trường \texttt{cv\_embedding} kiểu \texttt{vector(1536)} lưu vector embedding để phục vụ tìm kiếm ngữ nghĩa và matching.

\textbf{Bảng \texttt{cv\_drafts}} lưu bản nháp CV từ quá trình tải lên, gồm tệp gốc, văn bản đã trích xuất, và kết quả phân tích.

\textbf{Bảng \texttt{jd\_targets}} lưu các mô tả công việc mà người dùng muốn so khớp với CV. Mỗi JD cũng có \texttt{jd\_embedding} riêng để tránh tính lại.

\textbf{Bảng \texttt{match\_runs}} lưu lịch sử các lần so khớp -- mỗi dòng tương ứng với một (CV, JD, phương pháp), gồm điểm số, độ trễ, chi phí. Vừa làm cache vừa làm audit log cho thí nghiệm Chương 5.

\textbf{Bảng \texttt{audit\_logs}} ghi nhận các hành động của người dùng để theo dõi và phân tích.

\section{Thiết kế API}

Hệ thống cung cấp 19 API endpoint được tổ chức theo nhóm chức năng (Bảng \ref{tab:api}). Mọi API trả về JSON; các API yêu cầu xác thực kiểm tra JWT được lưu trong cookie httpOnly.

\begin{table}[H]
\centering
\caption{Danh sách các API endpoint chính}
\label{tab:api}
\resizebox{\textwidth}{!}{%
\begin{tabular}{|l|l|p{6cm}|}
\hline
\textbf{Endpoint} & \textbf{Phương thức} & \textbf{Mô tả} \\ \hline
\multicolumn{3}{|c|}{\textbf{Xác thực}} \\ \hline
/api/register & POST & Đăng ký tài khoản mới \\ \hline
/api/login & POST & Đăng nhập email/mật khẩu \\ \hline
/api/auth/logout & POST & Đăng xuất \\ \hline
/api/auth/me & GET & Lấy thông tin người dùng hiện tại \\ \hline
/api/auth/google/\{signin,callback\} & GET & OAuth Google \\ \hline
/api/auth/linkedin/\{signin,callback\} & GET & OAuth LinkedIn \\ \hline
\multicolumn{3}{|c|}{\textbf{Tải lên \& Phân tích CV}} \\ \hline
/api/upload/cv & POST & Upload tệp PDF/DOCX và phân tích \\ \hline
/api/upload/cv-blob & POST & Upload tệp lên Vercel Blob \\ \hline
/api/cv/\{cvId\} & GET, POST & Lấy/cập nhật dữ liệu CV \\ \hline
/api/cv-parser-stats & GET & Thống kê hiệu năng parser \\ \hline
\multicolumn{3}{|c|}{\textbf{Matching CV--JD}} \\ \hline
/api/cv/match & POST & Chạy 3 phương pháp matching song song \\ \hline
/api/cv/embed & POST & Tính embedding cho CV hoặc JD \\ \hline
/api/cv/jd-analyze & POST & Phân tích chi tiết CV vs JD bằng LLM \\ \hline
/api/cv/rewrite-bullet & POST & Viết lại một mục bằng LLM \\ \hline
/api/cv/summary & POST & Sinh tóm tắt nghề nghiệp bằng LLM \\ \hline
\multicolumn{3}{|c|}{\textbf{Quản trị}} \\ \hline
/api/admin/create & POST & Tạo tài khoản admin \\ \hline
\end{tabular}%
}
\end{table}

\section{Thiết kế pipeline phân tích CV hybrid}

Pipeline phân tích CV là thành phần cốt lõi của hệ thống, sử dụng phương pháp lai kết hợp xử lý bằng biểu thức chính quy và LLM. Hình \ref{fig:pipeline} mô tả luồng xử lý tổng thể.

\begin{figure}[H]
\centering
\begin{tikzpicture}[
  node distance=0.7cm and 0.9cm,
  every node/.style={font=\small},
  pstep/.style={rectangle, rounded corners, draw, thick, minimum width=3.2cm, minimum height=0.9cm, align=center, fill=blue!8},
  parser/.style={rectangle, rounded corners, draw, thick, minimum width=3.2cm, minimum height=0.9cm, align=center, fill=orange!15},
  data/.style={cylinder, draw, thick, aspect=0.25, shape border rotate=90, minimum height=1cm, minimum width=2.2cm, align=center, fill=gray!10},
  arr/.style={-{Stealth[length=2.5mm]}, thick}
]
\node[pstep] (upload) {Tải lên CV\\(PDF / DOCX)};
\node[pstep, right=of upload] (extract) {Trích xuất văn bản};
\node[pstep, right=of extract] (lang) {Nhận diện\\ngôn ngữ};

\node[parser, below=1.3cm of extract, xshift=-2.0cm] (regex) {Phân tích bằng\\biểu thức chính quy\\(liên hệ, kỹ năng)};
\node[parser, below=1.3cm of extract, xshift=2.0cm] (llm) {Phân tích bằng LLM\\(kinh nghiệm\\làm việc)};

\node[pstep, below=1.3cm of regex, xshift=2.0cm, fill=green!12] (merge) {Hợp nhất kết quả\\+ tính điểm tin cậy};
\node[data, below=0.7cm of merge] (db) {Cơ sở dữ liệu\\(cv\_workflow)};
\node[pstep, right=1.5cm of db] (editor) {Trình soạn thảo\\CV};

\draw[arr] (upload) -- (extract);
\draw[arr] (extract) -- (lang);
\draw[arr] (lang.south) -| (llm.north);
\draw[arr] (lang.south) -| (regex.north);
\draw[arr] (regex.south) -- (merge.north -| regex.south);
\draw[arr] (llm.south) -- (merge.north -| llm.south);
\draw[arr] (merge) -- (db);
\draw[arr] (db) -- (editor);
\end{tikzpicture}
\caption{Pipeline phân tích CV hybrid (kết hợp regex và LLM)}
\label{fig:pipeline}
\end{figure}

Pipeline hoạt động theo các bước sau:

\textbf{Bước 1 -- Tải lên và trích xuất văn bản:} Người dùng tải tệp PDF hoặc DOCX. Hệ thống dùng \texttt{pdf-parse} cho PDF hoặc \texttt{mammoth} cho DOCX. Tệp gốc lưu vào Vercel Blob.

\textbf{Bước 2 -- Nhận diện ngôn ngữ:} Module nhận diện ngôn ngữ phân tích văn bản để xác định CV viết bằng tiếng Việt hay tiếng Anh dựa trên ký tự Unicode đặc trưng và tần suất từ vựng.

\textbf{Bước 3 -- Phân tích bằng regex:} Các trường có cấu trúc rõ được trích trực tiếp bằng biểu thức chính quy: email, số điện thoại, URL LinkedIn, danh sách kỹ năng dựa trên từ khoá phổ biến.

\textbf{Bước 4 -- Phân tích bằng LLM:} Phần kinh nghiệm làm việc (có cấu trúc phức tạp, viết tự do) được gửi tới ChatGPT API với prompt chuyên biệt. Prompt yêu cầu mô hình trả về JSON với các trường: vị trí, công ty, thời gian, mô tả công việc.

\textbf{Bước 5 -- Hợp nhất và tính điểm tin cậy:} Kết quả từ regex và LLM được gộp thành cấu trúc JSON thống nhất. Hệ thống tính \texttt{possibility\_score} (0--100) dựa trên mức độ đầy đủ của thông tin trích xuất.

\textbf{Bước 6 -- Lưu trữ và hiển thị:} Dữ liệu CV có cấu trúc được lưu vào bảng \texttt{cv\_workflow} kiểu JSONB và chuyển đến trình soạn thảo để người dùng xem và chỉnh sửa.

Các bước 3 (regex) và 4 (LLM) độc lập với nhau và đều phụ thuộc vào kết quả bước 2, vì vậy trong hiện thực được thực thi \textit{song song} qua \texttt{Promise.all} nhằm giảm độ trễ tổng thể.

\section{Thiết kế quy trình so sánh CV--JD}

Để trả lời hai câu hỏi nghiên cứu RQ1 và RQ2, hệ thống cài đặt ba phương pháp so khớp CV--JD và chạy song song trên cùng một cặp đầu vào. Hình \ref{fig:matching_pipeline} mô tả luồng xử lý.

\begin{figure}[H]
\centering
\begin{tikzpicture}[
  node distance=0.8cm and 1.2cm,
  every node/.style={font=\small},
  inp/.style={rectangle, rounded corners, draw, thick, minimum width=3.0cm, minimum height=0.8cm, align=center, fill=gray!15},
  proc/.style={rectangle, rounded corners, draw, thick, minimum width=3.5cm, minimum height=1.1cm, align=center},
  outbox/.style={rectangle, rounded corners, draw, thick, minimum width=3.5cm, minimum height=0.8cm, align=center, fill=green!10},
  arr/.style={-{Stealth[length=2.5mm]}, thick}
]
\node[inp] (cv) at (0,0) {Văn bản CV};
\node[inp, right=1.2cm of cv] (jd) {Văn bản JD};

\node[proc, fill=blue!10, below=1.5cm of cv, xshift=1.5cm] (tfidf) {TF-IDF Matcher\\(corpus IDF, cosine)};
\node[proc, fill=orange!12, below=1.5cm of tfidf] (embed) {Embedding Matcher\\(text-embedding-3-small)};
\node[proc, fill=purple!10, below=1.5cm of embed] (llm) {LLM Matcher\\(gpt-4o-mini, temp=0)};

\node[outbox, right=2cm of tfidf] (s1) {Điểm + matched\_keywords};
\node[outbox, right=2cm of embed] (s2) {Điểm + cosine sim};
\node[outbox, right=2cm of llm]   (s3) {Điểm + reasoning};

\draw[arr] (cv.south) -| ($(tfidf.north)+(-0.5,0)$);
\draw[arr] (jd.south) -| ($(tfidf.north)+(0.5,0)$);
\draw[arr] (cv.south) -| ($(embed.north)+(-0.5,0)$);
\draw[arr] (jd.south) -| ($(embed.north)+(0.5,0)$);
\draw[arr] (cv.south) -| ($(llm.north)+(-0.5,0)$);
\draw[arr] (jd.south) -| ($(llm.north)+(0.5,0)$);

\draw[arr] (tfidf.east) -- (s1.west);
\draw[arr] (embed.east) -- (s2.west);
\draw[arr] (llm.east) -- (s3.west);
\end{tikzpicture}
\caption{Pipeline so sánh CV--JD bằng ba phương pháp chạy song song}
\label{fig:matching_pipeline}
\end{figure}

Mỗi phương pháp trả về một bản ghi gồm: điểm chuẩn hoá 0--100, danh sách từ khoá khớp/thiếu (đối với TF-IDF và LLM), độ trễ (ms), chi phí (USD), số token đã dùng. Kết quả của cả ba phương pháp được lưu vào bảng \texttt{match\_runs} để dùng cho hiển thị và cho thí nghiệm so sánh ở Chương 5. Nhờ chạy song song qua \texttt{Promise.allSettled}, một phương pháp lỗi không ảnh hưởng đến hai phương pháp còn lại.

\end{document}
